\begin{Exercise}[label={ex:matrice}]
    Una matrice sparsa è una matrice in cui la maggior parte degli elementi è costituita da zeri. 
    Spesso una matrice sparsa viene rappresentata in formato COO (Coordinate List), 
    ovvero come una sequenza di triple (i, j, valore) che indicano che, nella posizione (i, j), 
    è presente un valore diverso da zero (valore $\neq$ 0). Si consideri un file binario denominato 
    "input6.bin", che contiene due interi rappresentanti il numero di righe e di colonne della matrice, 
    seguiti dalla matrice stessa rappresentata in modalità ordinaria (cioè come sequenza di r * c elementi,
    dove r è il numero di righe e c il numero di colonne), i cui elementi sono numeri
    in virgola mobile a doppia precisione.
    Si scriva un programma che legga tale file e salvi in un secondo file di testo denominato "coo6.txt" 
    la rappresentazione COO della matrice.
\end{Exercise}

\begin{Answer}[ref={ex:matrice}]
  	\lstinputlisting[style=cstylesmall]{src/matrice.c}
\end{Answer}

